\chapter{Generalized Measurement and Dilation}
\label{ch:povm}

Chapter~\ref{ch:born-rule} deliberately deferred one statement:
$p_i=\operatorname{tr}(\rho E_i)$ for general effects $E_i$, not necessarily
rank-one projectors. This chapter closes that gap, and shows it reconnects
directly to Chapter~\ref{ch:partial-trace}'s partial trace — the two
constructions are complementary directions through the same
system-plus-environment architecture.

\section{Effects and POVMs}

\begin{definition}[POVM]
\label{def:povm}
\index{POVM (definition)}
A \emph{positive operator-valued measure} (POVM) is a family of operators
$\{E_i\}$ with $0 \preceq E_i \preceq I$ and $\sum_i E_i = I$. The rule
\begin{equation}
p_i = \operatorname{tr}(\rho E_i)
\end{equation}
generalizes Chapter~\ref{ch:observables}'s sharp measurements: a sharp
measurement is the special case $E_i = P_i$ with $P_i^2=P_i$; a general
$E_i$ need not be a projector, and need not be diagonal in any single
orthonormal basis shared by all the other $E_j$.
\end{definition}

\section{Dilation: Every POVM Is a Sharp Measurement, Elsewhere}

\begin{theorem}[Naimark dilation, finite-dimensional]
\label{thm:naimark}
\index{Naimark dilation, finite-dimensional (theorem)}
For any POVM $\{E_i\}_{i=1}^n$ on $\mathcal H$, there is an ancilla Hilbert
space $\C^n$, a reference state $|a_0\rangle \in \C^n$, a unitary $U$ on
$\mathcal H \otimes \C^n$, and a sharp measurement $\{\Pi_i = I_{\mathcal H}
\otimes |i\rangle\langle i|\}$ on $\mathcal H\otimes\C^n$, such that for every
state $\rho$ on $\mathcal H$,
\begin{equation}
\operatorname{tr}(\rho E_i) = \operatorname{tr}\Big[U\big(\rho \otimes
|a_0\rangle\langle a_0|\big)U^\dagger\, \Pi_i\Big].
\end{equation}
\end{theorem}

\begin{proof}
Define $V:\mathcal H \to \mathcal H\otimes\C^n$ by
\begin{equation}
V|\psi\rangle = \sum_{i=1}^n \big(\sqrt{E_i}\,|\psi\rangle\big) \otimes
|i\rangle,
\end{equation}
using the unique positive square root of each $E_i\succeq0$. For any
$|\phi\rangle,|\psi\rangle$,
\begin{equation}
\langle V\phi|V\psi\rangle = \sum_i \langle\phi|\sqrt{E_i}\sqrt{E_i}|\psi
\rangle = \langle\phi|\Big(\sum_i E_i\Big)|\psi\rangle = \langle\phi|\psi
\rangle,
\end{equation}
using $\sum_i E_i = I$: $V$ is an isometry. Fix $|a_0\rangle := |1\rangle \in
\C^n$ (any basis vector). Since $V$ is an isometry, the assignment
$U(|\psi\rangle\otimes|a_0\rangle) := V|\psi\rangle$ is well-defined and
inner-product-preserving on the subspace $\mathcal H\otimes|a_0\rangle$; by
the standard extension theorem for isometries between finite-dimensional
Hilbert spaces (extend an orthonormal basis of the range to a full
orthonormal basis of the codomain), $U$ extends to a unitary on the whole
space $\mathcal H\otimes\C^n$.

For a pure state $|\psi\rangle$, a direct computation (expanding
$(I\otimes|i\rangle\langle i|)V|\psi\rangle = \big(\sqrt{E_i}|\psi\rangle
\big)\otimes|i\rangle$ and taking the inner product with $V|\psi\rangle$)
gives
\begin{equation}
\langle V\psi|\Pi_i|V\psi\rangle = \langle\psi|\sqrt{E_i}\sqrt{E_i}|\psi
\rangle = \langle\psi|E_i|\psi\rangle = \operatorname{tr}(|\psi\rangle\langle
\psi|\,E_i),
\end{equation}
and since $U(|\psi\rangle\otimes|a_0\rangle) = V|\psi\rangle$, this is exactly
$\operatorname{tr}\big[U(|\psi\rangle\langle\psi|\otimes|a_0\rangle\langle
a_0|)U^\dagger\,\Pi_i\big]$. Both sides of the claimed identity are linear in
$\rho$, so the pure-state case extends to general $\rho$ by writing $\rho =
\sum_k q_k|\psi_k\rangle\langle\psi_k|$ and summing.
\end{proof}

\begin{remark}
This makes precise the sense in which every generalized measurement is an
ordinary sharp measurement performed elsewhere: couple the system to an
ancilla prepared in a fixed reference state, evolve unitarily, and measure
projectively on the composite. Nothing about the POVM formalism is a new kind
of physical process; it is the statistics of Chapters
\ref{ch:born-rule}--\ref{ch:observables}, viewed from a subsystem of a larger
sharp measurement.
\end{remark}

\section{The Same Architecture as Partial Trace}

\begin{remark}[Two directions through one construction]
Chapter~\ref{ch:partial-trace} asked: given a joint state on
$\mathcal H\otimes\mathcal H'$, what does a local observer restricted to
$\mathcal H$ see? The partial trace answers this by discarding $\mathcal H'$
entirely. This chapter asks the complementary question: given only local
access to $\mathcal H$, what joint constructions on $\mathcal H\otimes
\mathcal H'$ could have produced a given set of local statistics? Theorem
\ref{thm:naimark} answers this by exhibiting $\mathcal H'=\C^n$ explicitly.
The two constructions are literally inverse in spirit — one forgets an
already-composite system's inaccessible part; the other manufactures the
composite that a given (possibly non-projective) local measurement scheme
could have come from — and both rest on exactly the same system-plus-ancilla
architecture introduced for entanglement in Chapter~\ref{ch:entanglement}.
\end{remark}

\begin{ledgerentry}[End of Chapter~\ref{ch:povm}]
\textbf{Established:} every finite POVM admits an explicit Naimark dilation
into a sharp measurement on a system-plus-ancilla composite
(Theorem~\ref{thm:naimark}, proved constructively, not merely cited); this
closes the last gap flagged in Chapter~\ref{ch:born-rule}'s ledger.\\
\textbf{Not established:} anything about POVMs with a continuum of outcomes
(this book's finite-dimensional scope stops at finite $n$); the operational
question of which unitary $U$ and which ancilla a given real-world apparatus
actually realizes is not addressed — Theorem~\ref{thm:naimark} guarantees
\emph{some} dilation exists, not that it is unique or physically
distinguished; instruments and post-measurement state updates (the analogue
of $\rho \mapsto \sqrt{E_i}\rho\sqrt{E_i}^\dagger/\operatorname{tr}(\rho E_i)$)
are not derived here.
\end{ledgerentry}
